\documentclass[11pt]{article}

\usepackage[a4paper,margin=1in]{geometry}
\usepackage{amsmath,amssymb}
\usepackage{hyperref}

\hypersetup{
    colorlinks=true,
    linkcolor=blue,
    citecolor=blue,
    urlcolor=blue
}

\title{The First Marginal of a Transport Plan}
\author{}
\date{}

\begin{document}

\maketitle

Let
\[
\gamma \in \mathcal P(\mathbb R^d\times \mathbb R^d)
\]
be a transport plan. The projection onto the first coordinate is the map
\[
\pi_x:\mathbb R^d\times\mathbb R^d\to\mathbb R^d,
\qquad
\pi_x(x,y)=x.
\]
The condition
\[
(\pi_x)_\#\gamma=\mu
\]
means that if we ignore the destination coordinate $y$ and keep only the source
coordinate $x$, then the resulting distribution is $\mu$.

Explicitly, for every Borel set $A\subseteq\mathbb R^d$,
\[
\mu(A)
=
(\pi_x)_\#\gamma(A)
=
\gamma(\pi_x^{-1}(A))
=
\gamma(A\times\mathbb R^d).
\]
Thus the total amount of mass in the plan whose source lies in $A$ must equal
the original mass $\mu(A)$.

\section*{Finite Example}

Suppose
\[
\mu=\frac12\delta_0+\frac12\delta_1
\]
and define a transport plan
\[
\gamma
=
\frac12\delta_{(0,2)}
+
\frac12\delta_{(1,3)}.
\]
Then
\[
(\pi_x)_\#\gamma
=
\frac12\delta_{\pi_x(0,2)}
+
\frac12\delta_{\pi_x(1,3)}
=
\frac12\delta_0+\frac12\delta_1
=
\mu.
\]
The first marginal condition says that the plan starts with exactly the source
distribution $\mu$.

\end{document}
